% !TeX encoding = UTF-8
\chapter{Introduction}\label{chap:introduction}

In the second lab of the Security Insider Lab, we detect network attacks with machine learning. Instead of inspecting packets one by one, we describe each network connection by a vector of statistical features (a \emph{flow}) and train classifiers to tell benign flows from attack flows. The lab has four tasks:

\begin{enumerate}\setlength\itemsep{0pt}
	\item \textbf{Task 1:} we capture our own traffic and convert it into flows with CICFlowMeter.
	\item \textbf{Task 2:} we analyze a labeled flow dataset, CIC-IDS2017.
	\item \textbf{Task 3:} we pre-process the dataset: we handle infinite and missing values, scale the features, and select the 20 most informative ones.
	\item \textbf{Task 4:} we train Logistic Regression and Decision Tree classifiers, evaluate them, and apply the best model to the flows of our own capture.
\end{enumerate}

\paragraph{Dataset.} As suggested in the lab slides, we use the Friday-afternoon DDoS file of CIC-IDS2017, \texttt{Friday-WorkingHours-Afternoon-DDos.csv}. It contains benign flows and flows of a distributed denial-of-service attack, so the task is a binary classification problem, BENIGN vs.\ DDoS.

\paragraph{Structure.} Chapter~\ref{chap:task1_flows} describes how we produced flows from our own capture and the problems we ran into. Chapter~\ref{chap:task2_dataset} introduces the dataset, Chapter~\ref{chap:task3_preprocessing} the pre-processing, and Chapter~\ref{chap:task4_ml} the classifiers and their evaluation. Chapter~\ref{chap:task4_realtime} applies the best model to our own flows, and Chapter~\ref{chap:discussion} discusses the results and limitations. All code is in our \texttt{ddosids} Python package; the listings in this report are excerpts of the real source files.
